\section{Introduction}\label{sec:introduction}

Equatorial Plasma Bubbles (EPBs) are large-scale ionospheric plasma density depletions that form in the post-sunset equatorial F-region through the Generalized Rayleigh-Taylor instability. As the solar terminator passes across the magnetic equator, rapid recombination in the lower ionosphere creates a steep vertical density gradient below the denser F-layer. When accompanied by a strong pre-reversal enhancement of the eastward electric field, low-density plasma plumes rise into the topside ionosphere and drift eastward across low latitudes. Because these depleted channels contain steep electron density irregularities, they cause amplitude and phase scintillations in radio signals, degrading Global Navigation Satellite Systems (GNSS) positioning and satellite communications \cite{zhong2025automatic,okoh2025bootstrapping,githio2024machine}.

Ground-based optical monitoring of EPBs typically relies on narrow-band ($630.0\text{ nm}$ atomic oxygen red-line) All-Sky Imagers (ASIs), which record nocturnal airglow emissions at altitudes of $250\text{--}300\text{ km}$ \cite{okoh2025bootstrapping,zhong2025automatic}. In routine observatory operations, identifying plasma bubbles requires visual inspection of hundreds of exposures per night. Manual analysis is slow and often ambiguous when faint depletions overlap with bright lunar glare, horizon brightening from the Van Rhijn effect, Milky Way bands, or intermittent tropospheric clouds.

Several recent works have applied machine learning and deep learning to automate airglow image analysis. Convolutional Neural Networks (CNNs) and segmentation models have been used to classify EPB occurrences, track medium-scale traveling ionospheric disturbances, and build multi-year climatologies \cite{srisamoodkham2022detecting,liu2022statistical,sedlak2023analysis,okoh2025bootstrapping,zhong2025automatic,zhong2025statistical}. Other studies have paired CNN feature extractors with Support Vector Machines (SVMs) or collocated GNSS scintillation sensors \cite{thanakulketsarat2023classification,thanakulketsarat2023equatorial,githio2024machine}. However, training deep neural networks on moderate-sized, station-specific optical datasets carries a high computational cost and a clear risk of shortcut learning, where models memorize background luminance levels, dome struts, or embedded timestamp text rather than plasma depletion morphology. To preserve two-dimensional spatial structure without the parameter overhead of deep networks, Yacoub et al. \cite{yacoub2025automatic} and Siddiqui \cite{siddiqui2025enhanced} recently combined Two-Dimensional Principal Component Analysis (2DPCA) \cite{yang2004two} with tree-based and ensemble classifiers for airglow imagery.

Despite these advances, three practical issues remain open in automated ASI monitoring. First, local contrast enhancement methods amplify dark cloud shadows into sharp bands that standard classifiers mistake for EPBs, making an upstream sky quality filter necessary before classification. Second, consecutive airglow exposures taken $3\text{--}5\text{ minutes}$ apart share strong temporal correlation, meaning that random frame-level splits can overestimate generalization unless validated across grouped time windows and unseen observatories. Third, raw all-sky images contain burned-in station codes and timestamps in their corners, which must be tested as potential confounders.

In this paper, we present a two-stage cascaded pipeline for automated EPB detection in $630.0\text{ nm}$ all-sky imagery from the Brazilian EMBRACE network. The first stage applies an astronomical sky quality control check based on point-source star counting and mean sky brightness within the circular camera dome to reject cloud-obscured exposures. The second stage applies Gaussian background subtraction and adaptive histogram equalization, projects each $64 \times 64$ image matrix onto a compact 2DPCA subspace, and classifies the resulting features using a stacking ensemble \cite{wolpert1992stacked} of SVM, XGBoost, and Random Forest learners with in-fold SMOTE balancing \cite{chawla2002smote}. We evaluate the pipeline across stratified, temporally grouped, FOV-masked, and zero-shot cross-station protocols, and we inspect its spatial decisions using both 2DPCA reconstruction residuals and pixel-level SHapley Additive exPlanations (SHAP) \cite{lundberg2017unified}.
